\chapter{Implementation}
\label{ch:implementation}

This chapter describes the software that realises the design: its structure
(\cref{sec:impl:structure}), the core components (\cref{sec:impl:core}), the attack simulators
(\cref{sec:impl:attacks}), the language-model weight generator (\cref{sec:impl:stub}), the test
suite (\cref{sec:impl:tests}), and three defects found during evaluation and how they were fixed
(\cref{sec:impl:bugs}).

% ---------------------------------------------------------------------------
\section{Structure and dependencies}
\label{sec:impl:structure}

\textsc{EdgeGuard} is a pure-Python package that depends on NumPy, SciPy, PyYAML and the
\code{cryptography} library. PyTorch is an optional dependency, needed only for the forward-hook
adapter. The repository layout is:

\begin{lstlisting}[language={},numbers=none,frame=single,basicstyle=\ttfamily\scriptsize]
edgeguard/
|-- src/edgeguard/
|   |-- attestation/     merkle.py, manager.py         Merkle tree, signed attestation
|   |-- detection/       statistics.py, calibration.py, monitor.py
|   |-- classifier/      classifier.py, rules.py        base threat classifier
|   |-- llm/             model.py, detector.py          LLM weight stub, LLM detector
|   |-- attacks/         bitflip.py, weight_extraction.py, training_data_recovery.py,
|   |                    supply_chain.py, adversarial.py, synergy.py
|   |-- instrumentation/ hooks.py, qnn_adapter.py      PyTorch hooks, QNN protocol
|   `-- utils/           crypto.py, quantization.py
|-- tests/               test_merkle.py, test_monitor.py, test_classifier.py,
|                        test_llm_detector.py, attacks/test_attacks.py
|-- experiments/         run_evaluation.py, make_figures.py, make_tables.py, results/
|-- scripts/             build_merkle.py, calibrate.py, evaluate.py   (CLI tools)
|-- configs/default.yaml
`-- report/              this document; figures/ and generated/ are produced by experiments/
\end{lstlisting}

\Cref{tab:modules} lists the modules with their sizes.

\begin{table}[t]
\centering
\caption{Source modules and line counts (including docstrings and comments).}
\label{tab:modules}
\small
\begin{tabular}{llr}
\toprule
\textbf{Package} & \textbf{Module} & \textbf{Lines} \\
\midrule
attestation     & \code{merkle.py} (tree, proofs, localisation)          & 222 \\
                & \code{manager.py} (build/verify/save/load, ECDSA)      & 232 \\
detection       & \code{statistics.py} (Welford, $z$, Mahalanobis)        & 231 \\
                & \code{calibration.py} (calibration driver, JSON I/O)   & 178 \\
                & \code{monitor.py} (per-inference scoring)              & 214 \\
classifier      & \code{classifier.py} (base classifier, explanations)   & 278 \\
                & \code{rules.py} (confidence heuristics, presets)       & 139 \\
llm             & \code{model.py} (INT8 weight stub, 5 profiles)          & 209 \\
                & \code{detector.py} (LLM scores, priority rules)        & 320 \\
attacks         & \code{bitflip.py}                                      & 217 \\
                & \code{weight\_extraction.py}                           & 204 \\
                & \code{training\_data\_recovery.py}                      & 259 \\
                & \code{supply\_chain.py}                                & 215 \\
                & \code{adversarial.py}                                  & 153 \\
                & \code{synergy.py}                                      & 173 \\
instrumentation & \code{hooks.py}, \code{qnn\_adapter.py}                & 380 \\
utils           & \code{crypto.py}, \code{quantization.py}               & 264 \\
\midrule
\multicolumn{2}{l}{Total, including \code{\_\_init\_\_.py} files} & $\approx$4{,}050 \\
\bottomrule
\end{tabular}
\end{table}

% ---------------------------------------------------------------------------
\section{Core components}
\label{sec:impl:core}

\subsection{Hashing and signing}

\Cref{lst:hash} shows the leaf hash of \cref{eq:leaf}. The shape is packed as big-endian
unsigned 64-bit integers before the raw tensor bytes. \code{tobytes()} returns C-order bytes
regardless of the array's memory layout, so a transposed view and a contiguous copy with the
same logical contents hash identically.

\begin{lstlisting}[caption={Leaf hashing (\code{utils/crypto.py}).},label={lst:hash}]
def hash_weight_tensor(weights: np.ndarray) -> bytes:
    h = hashlib.sha256()
    shape_bytes = struct.pack(f">{len(weights.shape)}Q", *weights.shape)
    h.update(shape_bytes)
    h.update(weights.tobytes())
    return h.digest()

def hash_pair(left: bytes, right: bytes) -> bytes:
    h = hashlib.sha256()
    h.update(left)
    h.update(right)
    return h.digest()
\end{lstlisting}

Signing uses ECDSA over P-256 with SHA-256 from the \code{cryptography} library. If no key is
supplied, \code{MerkleAttestationManager.build} generates an ephemeral key pair, which is how
the experiments run. The attestation, including the leaves, the root, the DER-encoded
signature and the PEM public key, serialises to JSON. The dtype is not part of the leaf: an
\code{int8} and a \code{uint8} tensor with identical bytes and shape hash identically. This is
harmless when the verifier controls how tensors are interpreted, but binding the dtype into the
leaf would be a cheap hardening step.

\subsection{Streaming statistics}

\Cref{lst:welford} shows the per-layer update. Activations are flattened and cast to
\code{float64} for accumulation. Integer accumulation would overflow quickly for the
second-moment terms, and the cost of the cast is negligible next to the outer product.

\begin{lstlisting}[caption={Welford update (\code{detection/statistics.py}), after the fix of \cref{sec:impl:bugs}.},label={lst:welford}]
def update(self, activation: np.ndarray) -> None:
    x = activation.flatten().astype(np.float64)
    ...
    self._n += 1
    delta = x - self._mean
    self._mean += delta / self._n
    delta2 = x - self._mean
    self._M2 += delta * delta2              # per-channel variance numerator
    if self._C is not None:                 # dense co-moment if d <= max_cov_dim
        self._C += np.outer(delta, delta2)
\end{lstlisting}

At finalisation, the variance is $M_2/(n-1)$, floored at $10^{-6}$ to avoid division by zero for
constant channels. The covariance is $C/(n-1)+10^{-4}I$; the small ridge keeps the
pseudo-inverse stable.

\subsection{Monitor and classifiers}

\code{LayerwiseActivationMonitor.inspect} loops over calibrated layers, computes $Z_\ell$ and
$D_\ell$, applies the two thresholds, and returns a \code{MonitorResult} with per-layer scores,
the anomalous set, input-level maxima and an explanation. The base \code{ThreatClassifier}
implements the weight-evidence and activation-evidence rules of \cref{sec:design:classifier}.
\code{LLMThreatDetector} wraps it, adds the memorisation and entropy scores, and applies the full
priority order of \cref{fig:classifier}. Calibration statistics serialise to versioned JSON, so
that a deployment can ship them alongside the attestation.

\subsection{Command-line tools}

Three scripts wrap the library for use outside Python. \code{build\_merkle.py} generates a key
pair, builds an attestation from an \code{.npz} archive of weights, and verifies one.
\code{calibrate.py} fits monitor statistics from an archive of activations. \code{evaluate.py} is
a three-scenario smoke test for the base classifier, retained from the initial prototype. The
experiments in \cref{ch:evaluation} use \code{experiments/run\_evaluation.py} instead.

% ---------------------------------------------------------------------------
\section{Attack simulators}
\label{sec:impl:attacks}

Each simulator is a class with a \code{seed} argument and returns a dataclass describing what it
did. Weight attacks copy the weight dictionary before modifying it, a property that the tests
check. \Cref{tab:simulators} summarises their parameters.

\begin{table}[t]
\centering
\caption{Attack simulators and the parameters used in the evaluation.}
\label{tab:simulators}
\small
\begin{tabularx}{\textwidth}{@{}p{2.8cm}p{2.5cm}X@{}}
\toprule
\textbf{Simulator} & \textbf{Variant} & \textbf{What it does} \\
\midrule
\code{BitFlipAttack} & random & one bit at a uniformly random position, bit and tensor \\
 & sign\_bit & bit 7 of a uniformly random weight in a tensor \\
 & targeted & bit 7 of the largest-magnitude weight in a named tensor \\
 & progressive & bit 7 of the weight whose flip gives the largest change over all tensors (a gradient-free stand-in for PBFA \cite{rakin2019bfa}) \\
\addlinespace
\code{SupplyChain-} & steganographic & writes a 16-bit trigger into the LSBs of the first 16 weights of the first tensor \\
\code{Poisoning} & targeted\_layer & overwrites one row of a named tensor with values in $[100,127)$ \\
 & distributed & in every tensor, adds an integer in $[-3,3]$ to $\max(1,\lfloor 10^{-4}|L_i|\rfloor)$ random weights \\
\addlinespace
\code{AdversarialInput-} & text\_injection & early layers $\mathcal{N}(20, 12.5^2)$; late layers $\mathcal{N}(5, 6^2)$ \\
\code{Attack} & safety\_bypass & early $\mathcal{N}(15, 10^2)$; late $\mathcal{N}(25, 2.5^2)$ \\
 & membership\_probe & early $\mathcal{N}(8, 5.5^2)$; late $\mathcal{N}(30, 1.5^2)$ (not used in the main evaluation) \\
\addlinespace
\code{TrainingData-} & probe & early layers as clean ($\mathcal{N}(10,5^2)$); late layers mean $15+40f$ where $f$ is the relative depth, std $1.5$ \\
\code{Recovery} & & \\
\addlinespace
\code{WeightExtraction-} & memory\_dump & exact copy \\
\code{Attack} & cold\_boot & single-bit decay in 0.2\% of bytes (two's-complement correct) \\
 & em\_side\_channel & $\pm1$ LSB error in 2.5\% of weights \\
\addlinespace
\code{SynergyAttack} & targeted + safety\_bypass & targeted BFA on a named tensor; adversarial activations blended toward clean ones with weight 0.65 in early layers and 0.325 in late layers \\
\bottomrule
\end{tabularx}
\end{table}

All activation-level simulators draw clean activations from $\mathcal{N}(10, 5^2)$ clipped to the
INT8 range, the same distribution used for calibration. The attack activations are Gaussian
with different means and standard deviations, set per layer group as in \cref{tab:simulators}
(the parameters are absolute values, not offsets from the clean mean). This is the root of the circularity noted in \cref{ch:introduction}:
the attacks are defined as departures from a distribution that the detector models exactly. The
synergy simulator implements ``suppression'' as a convex blend toward fresh clean samples. It
does not model how a real attacker would craft such an input, and its success against the
monitor is a consequence of the blend weight, not evidence that real synergy attacks are
feasible.

The distributed-poisoning delta range includes zero, so a sampled weight may be left unchanged.
The evaluation therefore computes ground truth for localisation from the actual byte difference
between clean and attacked tensors, not from the simulator's report.

% ---------------------------------------------------------------------------
\section{Language-model weight stub}
\label{sec:impl:stub}

\code{QuantizedTransformerStub} produces an ordered dictionary of INT8 tensors with the layout
of \cref{fig:transformer_layout}. By default, values are drawn from $\mathcal{N}(0,25^2)$ and
clipped, which roughly fills the INT8 range like a symmetric per-tensor quantization of a
Gaussian weight matrix. Generation is lazy and cached. For the scaling experiment, the generator
is replaced by one that draws INT8 integers directly. This avoids the large float64 temporaries
that a Gaussian draw would create (2\,GiB for the Llama-3.2-1B embedding matrix alone), and it
does not affect hashing cost, which depends only on the byte count.

% ---------------------------------------------------------------------------
\section{Test suite}
\label{sec:impl:tests}

The suite contains \NTests{} tests, listed in full in \cref{app:tests}, and runs in about 40
seconds. \Cref{tab:test_summary} summarises it.

\begin{table}[t]
\centering
\caption{Test suite by file.}
\label{tab:test_summary}
\small
\begin{tabular}{lrp{8.2cm}}
\toprule
\textbf{File} & \textbf{Tests} & \textbf{Covers} \\
\midrule
\code{test\_merkle.py}         & 17 & tree shapes (1 leaf, powers of two, odd counts), determinism, single- and multi-layer tamper localisation, serialisation round trip, signed build/verify, save/load \\
\code{test\_monitor.py}        & 17 & Welford mean and variance accuracy, covariance against \code{np.cov} (regression), covariance skipped above $d_{\max}$, scoring functions, calibration save/load, threshold tuning \\
\code{test\_classifier.py}     & 15 & each base label, priority of weight evidence over activation evidence, rule presets and confidence functions \\
\code{test\_llm\_detector.py}  & 25 & stub structure and tensor counts for two profiles, dtype and range, every reachable LLM label path, score ranges \\
\code{attacks/test\_attacks.py}& 43 & each simulator's contract (copies input, determinism, effect size), two's-complement correctness (regression), five end-to-end attack-to-detection chains \\
\bottomrule
\end{tabular}
\end{table}

Several tests in \code{test\_llm\_detector.py} accept any of several labels, for example ``not
CLEAN'', rather than a single expected one. This was a deliberate choice during development,
because the label boundaries were still moving. It weakens those tests as specifications, and the
measured label behaviour in \cref{sec:eval:labels} is the more informative record.

% ---------------------------------------------------------------------------
\section{Defects found during evaluation}
\label{sec:impl:bugs}

Building the evaluation pipeline exposed three defects in code that had been passing its tests.
All three are fixed, each with a regression test, and all results in this report were produced
after the fixes.

\paragraph{1. Sign-bit overflow in the bit-flip simulator.} The original code computed a flip as
\code{np.int8(int(w[idx]) \^{} (1 << bit))}. For a negative weight and \code{bit} $=7$, the
Python integer XOR yields a value outside $[-128,127]$ (\cref{sec:bg:twos}), and NumPy 2 raises
\code{OverflowError}. The random strategy failed this way only when it happened to draw bit 7
for a negative weight, so short test runs could pass by chance. The fix XORs the
\code{uint8} view of the byte:

\begin{lstlisting}
w[idx] = (w[idx].view(np.uint8) ^ np.uint8(1 << bit)).view(np.int8)
\end{lstlisting}

The sign-bit, targeted and progressive strategies had the same pattern and received the same
fix.

\paragraph{2. Saturating instead of wrapping in the cold-boot model.} The cold-boot channel XORed
signed integers and then clipped the result to the INT8 range. A simulated decay of the sign bit
of $-5$ therefore produced $-128$ instead of $123$. The flip was silently replaced by a
saturation, which understates the damage of sign-bit decay. The fix again operates on the
unsigned view, vectorised over the selected bytes. A regression test checks that every corrupted
byte differs from the original in exactly one bit.

\paragraph{3. Double correction in the streaming covariance.} The co-moment update was written as
$C \mathrel{+}= \frac{n-1}{n}\,\boldsymbol\delta(\mathbf{x}-\boldsymbol\mu_n)^\top$. The
standard update is either $\boldsymbol\delta(\mathbf{x}-\boldsymbol\mu_n)^\top$ or
$\frac{n-1}{n}\boldsymbol\delta\boldsymbol\delta^\top$; the code combined both corrections.
Since $\mathbf{x}-\boldsymbol\mu_n = \frac{n-1}{n}\boldsymbol\delta$, each term was scaled by
$\big(\frac{n-1}{n}\big)^2$ instead of $\frac{n-1}{n}$, biasing the covariance low. Measured
against \code{np.cov}, the maximum relative error was 15.6\% after 20 samples and 1.2\% after
500. Per-channel variances, which the $z$-score uses, were computed separately and were correct.
The original tests checked the mean and variance but not the covariance, which is why the defect
survived. After the fix, every detection decision in the evaluation was unchanged. This indicates
that, with $\tau_D=50$, the $z$-score threshold made all the decisions in these experiments.

These defects share a lesson that matters for anyone building a bit-flip or statistics testbed.
Integer-width semantics and estimator algebra are easy to get subtly wrong in high-level code,
and tests written to check intended behaviour, rather than tests against an independent
reference such as \code{np.cov} or a bit-level popcount, can let such errors through.
